\chapter{Mercado Objetivo}

\section{Identificación del Mercado Objetivo}

\section{Identificación del Mercado Objetivo}

El crecimiento acelerado de las redes sociales y plataformas digitales ha generado una necesidad cada vez mayor de herramientas que fomenten entornos de interacción seguros y respetuosos. Si bien las grandes plataformas sociales ya cuentan con infraestructuras avanzadas de moderación, tanto automatizada como humana, esto no significa que el problema del lenguaje ofensivo esté resuelto. Sin embargo, la adopción de nuevas soluciones en estos entornos puede resultar limitada por las altas barreras tecnológicas y políticas.

Por ello, resulta más estratégico enfocar los esfuerzos iniciales en sectores donde la problemática persiste y las soluciones actuales son aún incipientes o poco accesibles. A continuación, se describen los mercados en los que la propuesta tecnológica tendría mayor impacto y viabilidad de implementación:

\begin{enumerate}
    \item \textbf{Plataformas de streaming en vivo:} Espacios como Kick o Twitch presentan una alta velocidad de interacción, con miles de mensajes por minuto durante transmisiones en directo. La moderación manual resulta insuficiente, por lo que un sistema automatizado puede actuar como apoyo eficaz, especialmente en contenidos dirigidos a públicos jóvenes o menores de edad.

    \item \textbf{Comunidades en línea:} Entornos con alta participación textual y fuerte presencia de usuarios hispanohablantes, donde la moderación suele depender únicamente de configuraciones básicas o listas de palabras prohibidas. Este grupo incluye servidores de Discord, grupos de Telegram y foros como Reddit. Cabe destacar que Reddit ha incorporado moderación automatizada mediante Perspective API, desarrollada por Google \cite{perspective2025research}.

    \item \textbf{Medios de comunicación digital:} Portales de noticias y sitios informativos que permiten la participación del público a través de comentarios. Estos espacios requieren mantener una conversación pública constructiva sin comprometer su reputación institucional. Medios como \textit{The New York Times} y \textit{El País} ya han incorporado tecnologías como Perspective API para gestionar este desafío \cite{perspective2025research}.

    \item \textbf{Plataformas educativas virtuales:} Aunque menos explotado, el sector educativo presenta un alto potencial, en especial en entornos con participación estudiantil activa. Como ejemplo destacado, el Departamento de Educación de Australia Meridional, en colaboración con Microsoft, ha
\end{enumerate}







\vspace{0.5cm}

implementado la herramienta ``EdChat'', apoyada por el servicio Azure AI Content Safety, con el objetivo de crear espacios digitales seguros para el aprendizaje.

\section{Elección del Mercado Objetivo}

Los medios digitales que permiten comentarios abiertos al público enfrentan un desafío constante en la gestión de expresiones ofensivas, discursos de odio y comportamientos disruptivos. Este tipo de contenido no solo perjudica la experiencia del lector, sino que también compromete la imagen institucional del medio y su capacidad para ofrecer espacios de diálogo respetuoso.

La elección de este sector responde a múltiples factores. En primer lugar, existe una necesidad evidente y persistente de soluciones eficientes, ya que muchas plataformas aún dependen de procesos manuales o herramientas básicas para moderar los comentarios. En segundo lugar, el volumen diario de interacciones es considerable, dificultando la revisión exhaustiva sin el apoyo de sistemas automatizados. Finalmente, el lenguaje ofensivo puede deteriorar gravemente la percepción pública del medio, afectando su credibilidad, su relación con los usuarios y su sostenibilidad a largo plazo.

La adopción exitosa de tecnologías como Perspective API evidencia esta demanda. De acuerdo con datos de Jigsaw (subsidiaria de Google), esta API procesa más de 500 millones de solicitudes diarias, y ha sido integrada en medios internacionales como \textit{The New York Times}, \textit{El País} y Cofina. Estos casos muestran cómo una herramienta de moderación automatizada puede reducir significativamente los niveles de toxicidad, mejorar la eficiencia del proceso editorial y ampliar la participación del público, al aumentar el número de artículos abiertos a comentarios y reducir la carga de la moderación humana \cite{jigsaw2021api}.

\section{Panorama Nacional en Medios Digitales}

En el contexto nacional, numerosos medios digitales han optado por restringir o suprimir los espacios públicos de comentarios en sus sitios web. Esta decisión suele estar motivada por la complejidad que implica la moderación efectiva de contenidos ofensivos, inapropiados o potencialmente dañinos. La ausencia de herramientas automatizadas adaptadas al idioma español acentúa esta dificultad, lo que lleva a muchas redacciones a limitar la interacción directa con su audiencia como medida preventiva.

Una situación comparable se presenta en programas de televisión que incluyen la participación del público en tiempo real, como transmisiones en vivo o formatos interactivos. En estos entornos, los mensajes enviados deben pasar por filtros previos a su difusión, con el objetivo de evitar expresiones discriminatorias, agresivas o inadecuadas para una audiencia general.




